\documentclass[10pt,a4paper]{amsart}
\usepackage[margin=19mm]{geometry}
\usepackage{amsmath,amssymb,mathtools}
\usepackage[scr=rsfs,cal=euler]{mathalfa}
\usepackage{xfrac}
\usepackage[unicode,hidelinks,bookmarks=false]{hyperref}
\newcommand{\ie}{i.e.\ }
\newcommand{\cf}{cf.\ }
\newcommand{\resp}{respectively\ }
\newcommand{\Subsection}{\subsection}
\providecommand{\nobreakdash}{\nobreak\hbox{-}}
\setlength{\emergencystretch}{2em}
\multlinegap0pt
\usepackage{bm}
\usepackage{bbm}
\usepackage{dsfont}
\usepackage{xspace}
\usepackage{cancel}
\usepackage{mathtools}
\usepackage{amsthm}
\makeatletter
\@ifundefined{thrm}{\newtheorem{thrm}{Thrm}}{}
\@ifundefined{dfn}{\newtheorem{dfn}[thrm]{Definition}}{}
\@ifundefined{lemma}{\newtheorem{lemma}[thrm]{Lemma}}{}
\makeatother
\makeatletter
\expandafter\ifx\csname frcseries\endcsname\relax
\newenvironment{frcseries}{\fontfamily{frc}\selectfont}{}
\fi
\def\p{{\varphi}}
\def\ph{{\Phi}}
\def\sb {{\nabla}}
\makeatother
\title[The Riemannian curvature identities of a $G\_2$ connection wi]{The Riemannian curvature identities of a $G\_2$ connection with skew-symmetric torsion and generalized Ricci solitons}
\author{Stefan Ivanov, Nikola Stanchev}
\date{Source: arXiv:2307.05619v5 (CC BY 4.0)}
\begin{document}
\maketitle

\noindent\emph{Source and licence.} The Riemannian curvature identities of a $G_2$ connection with skew-symmetric torsion and generalized Ricci solitons. Version arXiv:2307.05619v5 (CC BY 4.0), distributed under the Creative Commons Attribution 4.0 International licence, as recorded in the arXiv licence metadata for this exact version and shown on the abstract page https://arxiv.org/abs/2307.05619v5. Full rights notice: licenses/arxiv-2307.05619-CC-BY-4.0.txt.\par\smallskip

\noindent\textit{Editorial scope.} Counted unit: the Lemma whose source label is \texttt{None} in the article (source0.tex lines 965--967, source label \texttt{None}), delivered together with the article's own complete original proof of it (source0.tex lines 968--997, one \texttt{proof} environment of 30 lines). No mathematics is written, repaired or re-derived: statement and proof are the article's own text line for line. Required context retained uncounted: gen (lines 370--375); rri (lines 654--659); lmt (lines 733--735); nnew1 (lines 791--793); dtt (lines 830--838). Every definition, display and label of those lines is retained unchanged. Omitted from this package and not redistributed: 1--369, 376--653, 660--732, 736--790, 794--829, 839--964, 998--1668. Nothing inside a retained excerpt is omitted or altered: every retained body is delivered exactly as the article's lines are written, so no omission receipt of any kind accompanies this package. Citations: the retained excerpts carry no citation command at all, so no bibliography entry of the article is transcribed and no reference is redistributed. Reference adaptation: none was needed and none was made; every reference inside the delivered unit resolves inside it, so no unresolved reference or citation is produced. Printed slips kept literal: no printed word, symbol, hypothesis or number of the counted statement or of its proof was altered. Layout and class: the article's own class is \texttt{article} with packages including amsfonts, latexsym, amsmath, amssymb, amsthm, mathrsfs, hyperref; the wrapper is the lane's pinned \texttt{amsart} editorial wrapper at 10pt on A4, which restates the theorem-like environments the retained text needs and adds the article's own macro definitions that the retained text needs (\texttt{\textbackslash p}, \texttt{\textbackslash ph}, \texttt{\textbackslash sb}), with extra wrapper packages (bm, bbm, dsfont, xspace, cancel, mathtools). The delivered numbering is the unit's own. Assets: no retained span contains a figure environment, a graphics-inclusion command, a PDF-page inclusion, a table or a TikZ picture; no image file is copied into this package, the package declares no asset dependency and no image file is redistributed with it. Licence: version arXiv:2307.05619v5 of this article is distributed under the Creative Commons Attribution 4.0 International licence, as recorded in the arXiv licence metadata for this exact version and shown on the abstract page https://arxiv.org/abs/2307.05619v5. A full rights notice is delivered at licenses/arxiv-2307.05619-CC-BY-4.0.txt. Characters: every retained excerpt is pure ASCII, so the delivered engine, XeLaTeX with its default Latin Modern text font, reports no missing character.
 \begin{equation}\label{gen}
 \begin{split}
 R(X,Y,Z,V)+ R(Y,Z,X,V)+ R(Z,X,Y,V)-R(V,X,Y,Z)-R(V,Y,Z,X)-R(V,Z,X,Y)\\
 =\frac32dT(X,Y,Z,V)-\sigma^T(X,Y,Z,V).
 \end{split}
 \end{equation}

\begin{dfn}
The curvature $R$ of a linear connection on a $G_2$ manifold  is a $G_2$-instanton if  the curvature 2-form lies in the Lie algebra $\frak{g}_2\cong\Lambda^2_{14}$. This is equivalent to  the  identities:
\begin{equation}\label{rri}
R_{abij}\p_{abk}=0 \Longleftrightarrow R_{abij}\ph_{abkl}=-2R_{klij}.
\end{equation}
\end{dfn}

\begin{equation}\label{lmt}
T_{klm}\p_{klm}=-\frac32T_{jsk}\Phi_{jslm}\p_{klm}+42\lambda=-6T_{jsk}\p_{jsk}+42\lambda \Longrightarrow \lambda=\frac16T_{klm}\p_{klm}
\end{equation}

\begin{equation}\label{nnew1}
 -\delta T_{bc}\p_{bci}=\sb_sT_{sbc}\p_{bci}=\theta_sT_{skt}\p_{kti}+6d\lambda_i.
\end{equation}

\begin{equation}\label{dtt}
\begin{split}
dT_{iabc}\p_{abc}+2\sb_iT_{abc}\p_{abc}=
dT_{iabc}\p_{abc}+12d\lambda_i=0;\\
 3\sb_aT_{bci}\p_{abc}-2\sigma^T_{iabc}\p_{abc}-18d\lambda_i=0;\\
\sb_aT_{bci}\p_{abc}-2\theta_sT_{skt}\p_{kti}-6d\lambda_i=0.
%{\color{red}0=\Big[\sb_iT_{abc}+\sb_aT_{bic}+\sb_bT_{iac}-\sb_cT_{iab}+2\sigma^T_{iabc}\Big]\p_{abc}+12d\lambda_i}\\{\color{red}=\Big[-3\sb_aT_{bci}+2\sigma^T_{iabc}\Big]\p_{abc}+18d\lambda_i}
\end{split}
\end{equation}

\begin{lemma}
Let $(M,\p)$ be  an integrable $G_2$ manifold and the  curvature of the characteristic  connection $\sb$   is a $G_2$-instanton. Then  $\delta T\in\Lambda^2_{14}\cong  \frak{g}_2$.% the $G_2$ structure is of constant type and the torsion is closed, $d\lambda=dT=0$.
\end{lemma}

\begin{proof}
Suppose the curvature $R$ of $\sb$ is a $G_2$-instanton. % i.e. the curvature 2-form lies in the lie algebra $g_2\cong\Lambda^2_{14}$. 
%Taking into account \eqref{rr} and \eqref{rri}, we calculate applying the first Bianchi identity \eqref{1bi} that
%\begin{multline}\label{rri1}
%2Ric_{ji}=R_{abci}\ph_{abcj}=\frac13\Big[R_{abci}+R_{bcai}+R_{cabi} \Big]\ph_{abcj}=\frac13\Big[dT_{abci}-\sigma^T_{abci}+\sb_iT_{abc}  \Big]\ph_{abcj}\\
%0=R_{abci}\p_{abc}=\frac13\Big[R_{abci}+R_{bcai}+R_{cabi} \Big]\p_{abc}=\frac13\Big[dT_{abci}-\sigma^T_{abci}+\sb_iT_{abc}  \Big]\p_{abc}
%\end{multline}
Multiply \eqref{gen} with  $\p$ and  apply  \eqref{rri} to get
\begin{equation}\label{rri2}
0=\Big[3R_{abci}-3R_{iabc}\Big]\p_{abc}=\Big[\frac32dT_{abci}-\sigma^T_{abci}\Big]\p_{abc}
\end{equation}
We obtain from \eqref{rri2} and  \eqref{dtt} 
$$%\begin{equation}\label{rri3}
dT_{abci}\p_{abc}=\frac23\sigma^T_{abci}\p_{abc}=2\sb_iT_{abc}\p_{abc}
$$ %\end{equation} 
and use \eqref{lmt} to conclude 
\begin{equation}\label{rri5}
\begin{split}
\sb_iT_{abc}\p_{abc}=6d\lambda_i=\frac13\sigma^T_{abci}\p_{abc}=-\theta_s\p_{sab}T_{abi}=-\theta_sT_{sab}\p_{abi}=d^{\sb}\theta_{ab}\p_{abi}%;\\
%{\color{red} d^{\sb}\theta_{ab}\p_{abi}\p_{icd}(\ph_{abcd}+\delta_{ac}\delta_{bd}-\delta_{ad}\delta_{bc}}
\end{split}
\end{equation}
Substitute \eqref{rri5} into \eqref{dtt} and \eqref{nnew1} to get
\begin{equation}\label{rri6}
\begin{split}
\delta T_{ab}\p_{abi}=0 \Leftrightarrow \delta T\in\Lambda^2_{14};\quad \sb_aT_{bci}\p_{abc}=-6d\lambda_i.
\end{split}
\end{equation}
The lemma is proved.
\end{proof}

\end{document}
